\documentclass[10pt,a4paper]{article}
\usepackage{amsmath,amssymb,amsthm,graphicx}
\usepackage{amsmath,amsfonts,amssymb,amsthm,enumerate,hyperref,ytableau,authblk,listings}
\usepackage[capitalize]{cleveref}
\newtheorem{theorem}{Theorem}[section]
\newtheorem{lemma}[theorem]{Lemma}
\newtheorem{corollary}[theorem]{Corollary}
\newtheorem{conjecture}[theorem]{Conjecture}

\theoremstyle{definition}

\newtheorem{definition}[theorem]{Definition}

\providecommand{\coset}[2]{[#1 \mapsto #2]}
\providecommand{\dcoset}[2]{[#1 \leftrightarrow #2]}
\providecommand{\id}{\mathsf{id}}
\providecommand{\cM}{\mathcal{M}}
\title{Strict degree reduction for a Boolean function supported on one permutation coset}
\author{Gilad Chase, Neta Dafni, Yuval Filmus, Nathan Lindzey}
\date{arxiv:2210.00245v2; CC BY 4.0}
\begin{document}
\maketitle
\noindent\textit{Editorial note.} The spectral degree characterization and Young-antisymmetrizer spanning theorem are explicit allowed inputs, with their exact author statements supplied. Strictness under coset restriction is not assumed.\par\medskip
\noindent\textit{Editorial note.} The author credits this lemma to the proof of Dafni et al., Lemma5.7. That credit is retained. The finite computer searches elsewhere in the article are not used by this task.\par\medskip
\noindent\textit{Editorial scope.} One original degree-reduction lemma with its full proof. Standard symmetric-group decomposition, degree characterization, and tableau spanning are declared inputs. No2-intersecting-family extremal theorem, computational finite-case verification, or perfect-matching analogue is counted.\par\medskip
\noindent\textit{Allowed standard prerequisites (editorial).}\begin{itemize}
\item Isotypic decomposition and polynomial-degree characterization on S\_n. Real functions on S\_n decompose orthogonally into V\textasciicircum{}lambda. Their coordinate polynomial degree is the largest n-lambda\_1 among nonzero components.. source172-184, Ellis-Friedgut-Pilpel Theorem7.. Locate a highest-degree nonzero component of the restriction and infer the degree lower bound from the lifted component. Independent review must accept this as a stated standard representation-theoretic input, not an unproved new local lemma.
\item Young column-antisymmetrizer spanning theorem. V\textasciicircum{}lambda is spanned by chi\_(s,t)=sum over pi in C(t) of sign(pi)*e\_(s,t\textasciicircum{}pi), with the tableaux and e-functions defined in the source.. source331-364; Sagan, The Symmetric Group, Chapter2.. Choose a witness with nonzero inner product and lift it by adding a singleton row. The lift/cancellation argument is the counted proof obligation.
\end{itemize}\medskip
For integer $n \geq 1$, the \emph{symmetric group} $S_n$ is the collection of all permutations of $[n] = \{1,\ldots,n\}$. We denote the identity permutation by $\id$.

A \emph{$t$-coset} of $S_n$ is a set of the form
\[
 \{ \alpha \in S_n : \alpha(i_1) = j_1, \ldots, \alpha(i_t) = j_t \},
\]
where $i_1,\ldots,i_t \in [n]$ are distinct and $j_1,\ldots,j_t \in [n]$ are distinct. A $t$-coset contains $(n-t)!$ permutations.

We use the term \emph{coset} for a $1$-coset. For $i,j \in [n]$, the coset $\coset{i}{j}$ consists of all permutations sending $i$ to $j$:
\[
 \coset{i}{j} = \{ \alpha \in S_n : \alpha(i) = j \}.
\]

\subsubsection{Degree} \label{sec:degree-sym}

We can represent permutations $\alpha \in S_n$ using $n^2$ variables $x_{ij}$ whose semantics are: $x_{ij} = 1$ if $\alpha(i) = j$, and $x_{ij} = 0$ otherwise. The \emph{degree} of a function $f\colon S_n \to \mathbb{R}$, denoted $\deg f$, is the minimal $d$ such that $f$ can be represented as a degree~$d$ polynomial over the variables $x_{ij}$.
For example, the characteristic function of a $t$-coset has degree at most~$t$, since it can be represented by the polynomial $x_{i_1j_1} \cdots x_{i_tj_t}$.

An equivalent way to define degree is via the representation theory of the symmetric group. Let $\mathbb{R}[S_n]$ be the vector space of all real-valued functions over the symmetric group. Representation theory gives an orthogonal decomposition (with respect to the inner product $\langle f,g \rangle = \sum_{\alpha \in S_n} f(\alpha) g(\alpha)$)
\[
 \mathbb{R}[S_n] = \bigoplus_{\lambda \vdash n} V^\lambda,
\]
where $\lambda$ goes over all integer partitions of $n$, and $V^\lambda$ are certain subspaces known as \emph{isotypic components}, which we define explicitly in \Cref{sec:degree-reduction-sym}. Accordingly, every function $f\colon S_n \to \mathbb{R}$ has a unique decomposition
\[
 f = \sum_{\lambda \vdash n} f^{=\lambda},
\]
where $f^{=\lambda} \in V^\lambda$. 
Ellis, Friedgut and Pilpel~\cite[Theorem 7]{EFP} showed that $\deg f$ is the maximal $d$ such that $f^{=\lambda} \neq 0$ for some $\lambda$ satisfying $\lambda_1 = n - d$. In particular, every function has degree at most $n-1$.

\subsubsection{Degree reduction} \label{sec:degree-reduction-sym}

Let $\mathcal{F} \subseteq S_n$. We denote the restriction of $\mathcal{F}$ to the coset $\coset{i}{j}$ by $\mathcal{F}|_{\coset{i}{j}}$. We can think of $\mathcal{F}|_{\coset{i}{j}}$ as a subset of $S_{n-1}$. Formally speaking, we can think of $S_n$ as the set of all injections from $[n]$ to $[n]$. The restriction of $S_n$ to the coset $\coset{i}{j}$ is isomorphic to the set of all injections from $[n] \setminus \{i\}$ to $[n] \setminus \{j\}$, which is the same as $S_{n-1}$ up to renumbering.

Let $f$ be the characteristic function of $\mathcal{F}$. The definition of degree using polynomials shows that $\deg f|_{\coset{i}{j}} \leq \deg f$, where $f$ is the characteristic function of $\mathcal{F}|_{\coset{i}{j}}$. Indeed, given a polynomial representing $f$, we can obtain a polynomial representing $f|_{\coset{i}{j}}$ by substituting $x_{ij} = 1$ and $x_{ij'} = x_{i'j} = 0$ for any $i' \neq i$ and $j' \neq j$.

The following crucial lemma, which forms part of the proof of~\cite[Lemma 5.7]{DFLLV}, states that if $\mathcal{F}$ is contained in $\coset{i}{j}$, then the degree of $f$ strictly decreases when restricting to $\coset{i}{j}$.

The argument proving \Cref{lem:degree-reduction-sym} is more subtle. The starting point is the following explicit description of the subspaces $V^\lambda$, which we mentioned in \Cref{sec:degree-sym}.

Let $\lambda = \lambda_1,\ldots,\lambda_m \vdash n$. A \emph{Young tableau} of \emph{shape} $\lambda$ is an arrangement of the numbers $1,\ldots,n$ in $m$ rows of lengths $\lambda_1,\ldots,\lambda_m$, left justified. Here are some examples:
\[
 \begin{array}{ccc}
 \ytableaushort{341,2} && \ytableaushort{14,23} \\\\
 \lambda = (3,1) && \lambda = (2,2)
 \end{array}
\]

Let $s,t$ be a pair of Young tableaux of the same shape $\lambda_1,\ldots,\lambda_m \vdash n$. The function $e_{s,t}\colon S_n \to \{0,1\}$ is defined as follows: $e_{s,t}(\alpha) = 1$ if for every $k \in [m]$, $\alpha$ sends every number on the $k$-th row of $s$ to a number on the $k$-th row of $t$, and $e_{s,t}(\alpha) = 0$ otherwise. For example, consider
\[
 s = \ytableaushort{123,45} \quad
 t = \ytableaushort{315,24}
\]
Then $e_{s,t}(\alpha) = 1$ if $\alpha(1),\alpha(2),\alpha(3) \in \{3,1,5\}$ and $\alpha(4),\alpha(5) \in \{2,4\}$.

For a tableau $t$ with $m'$ columns, we denote by $C(t)$ the \emph{column stabilizer} of $t$, which is the set of all permutations of the entries of $t$ which only move entries within the same column. Each permutation $\pi \in C(t)$ thus consists of $m'$ permutations $\pi_1,\ldots,\pi_{m'}$, one for each column. The \emph{sign} of $\pi$, denote $(-1)^\pi$, is the product of the signs of the permutations $\pi_1,\ldots,\pi_{m'}$. We denote the result of applying $\pi$ to $t$ by $t^\pi$.

Let $s,t$ be a pair of Young tableaux of the same shape $\lambda \vdash n$. We define a function $\chi_{s,t}\colon S_n \to \{-1,0,1\}$ as follows:
\[
 \chi_{s,t} = \sum_{\pi \in C(t)} (-1)^\pi e_{s,t^\pi}.
\]
Continuing our above example, here is $t^\pi$ for all elements $\pi \in C(t)$, together with their sign:
\[
 \left(\ytableaushort{315,24}~,+1\right) \quad
 \left(\ytableaushort{345,21}~,-1\right) \quad
 \left(\ytableaushort{215,34}~,-1\right) \quad
 \left(\ytableaushort{245,31}~,+1\right)
\]

We are interested in the functions $\chi_{s,t}$ since they span $V^\lambda$.

\begin{theorem}[{\cite[Chapter 2]{Sagan}}] \label{thm:Vlambda-sym}
Let $\lambda \vdash n$. The subspace $V^\lambda$ is spanned by the functions $\chi_{s,t}$, where $s,t$ go over all pairs of Young tableaux of shape $\lambda$.
\end{theorem}

\begin{lemma} \label{lem:degree-reduction-sym}
Suppose that $f\colon S_n \to \{0,1\}$ is the characteristic function of a family which is a subset of the coset $\coset{i}{j}$. Then $\deg f|_{\coset{i}{j}} \leq \max( \deg f - 1, 0 )$.
\end{lemma}

\begin{proof}[Proof of \Cref{lem:degree-reduction-sym}]
Assume, for concreteness, that $i = j = n$.
If $n = 1$ then $\deg f|_{\coset{n}{n}} = 0$, so there is nothing to prove. If $f|_{\coset{n}{n}} = 0$ then again $\deg f|_{\coset{n}{n}} = 0$, so there is nothing to prove. Therefore we can assume that $n \geq 2$ and that $f|_{\coset{n}{n}} \neq 0$.

Let $d = \deg f|_{\coset{n}{n}}$. According to the spectral definition of degree (see \Cref{sec:degree-sym}), $f|_{\coset{n}{n}}^{=\lambda} \neq 0$ for some $\lambda \vdash n-1$ such that $\lambda_1 = (n-1) - d$. Since the decomposition into isotypic components is orthogonal, this implies that $\langle f|_{\coset{n}{n}}, f|_{\coset{n}{n}}^{=\lambda} \rangle = \langle f|_{\coset{n}{n}}^{=\lambda}, f|_{\coset{n}{n}}^{=\lambda} \rangle \neq 0$. Since $f|_{\coset{n}{n}}^{=\lambda} \in V^\lambda$, \Cref{thm:Vlambda-sym} implies that $\langle f|_{\coset{n}{n}}, \chi_{s,t} \rangle \neq 0$ for some Young tableaux $s,t$ of shape $\lambda$.

Let $s',t'$ be the Young tableaux obtained by adding a new row to $s,t$ (respectively) consisting only of the number $n$. Here is an example:
\[
 s = \ytableaushort{123,45} \quad
 t = \ytableaushort{315,24} \quad \longrightarrow \quad
 s' = \ytableaushort{123,45,6} \quad
 t' = \ytableaushort{315,24,6}
\]
The tableaux $s',t'$ have shape $\mu$ satisfying $\mu_1 = \lambda_1 = (n-1)-d = n-(d+1)$.
We will show that $\langle f, \chi_{s',t'} \rangle \neq 0$. The orthogonality of the decomposition into isotypic components would imply that $f^{=\mu} \neq 0$, and so $\deg f \geq d+1$ by the spectral definition of degree.

Let $\pi' \in C(t')$ be any permutation such that $\pi'(n) \neq n$. By assumption, $f(\alpha') = 1$ only if $\alpha'(n) = n$, in which case $e_{s',(t')^{\pi'}}(\alpha') = 0$. Therefore $\langle f, e_{s',(t')^{\pi'}} \rangle = 0$. This shows that
\[
 \langle f, \chi_{s',t'} \rangle = \sum_{\substack{\pi' \in C(t') \\ \pi'(n) = n}} (-1)^{\pi'} \langle f, e_{s',(t')^{\pi'}} \rangle.
\]
If $\pi'(n) = n$ then $e_{s',(t')^{\pi'}}(\alpha') \neq 0$ only if $\alpha'(n) = n$. Conversely, if $\alpha'(n) = n$, that is, if $\alpha' \in \coset{n}{n}$, then $e_{s',(t')^{\pi'}}(\alpha') = e_{s,t^\pi}(\alpha)$, where $\alpha,\pi$ are the restrictions of $\alpha',\pi'$ to $[n-1]$. Moreover, $(-1)^\pi = (-1)^{\pi'}$. Therefore
\[
 \langle f, \chi_{s',t'} \rangle = \sum_{\pi \in C(t)} (-1)^\pi \langle f|_{\coset{n}{n}}, e_{s,t^\pi} \rangle = \langle f|_{\coset{n}{n}}, \chi_{s,t} \rangle \neq 0. \qedhere
\]
\end{proof}

\newcommand{\etalchar}[1]{$^{#1}$}
\begin{thebibliography}{BMRR21}
\bibitem[DFL{\etalchar{+}}21]{DFLLV}
Neta Dafni, Yuval Filmus, Noam Lifshitz, Nathan Lindzey, and Marc Vinyals.
\newblock Complexity measures on the symmetric group and beyond.
\newblock In {\em 12th {I}nnovations in {T}heoretical {C}omputer {S}cience
  {C}onference}, volume 185 of {\em LIPIcs. Leibniz Int. Proc. Inform.}, pages
  Art. No. 87, 5. Schloss Dagstuhl. Leibniz-Zent. Inform., Wadern, 2021.

\bibitem[EFP11]{EFP}
David Ellis, Ehud Friedgut, and Haran Pilpel.
\newblock Intersecting families of permutations.
\newblock {\em J. Amer. Math. Soc.}, 24(3):649--682, 2011.
\bibitem[Sag01]{Sagan}
Bruce~E. Sagan.
\newblock {\em The symmetric group}, volume 203 of {\em Graduate Texts in
  Mathematics}.
\newblock Springer-Verlag, New York, second edition, 2001.
\newblock Representations, combinatorial algorithms, and symmetric functions.
\end{thebibliography}
\end{document}
